% !TeX root = ../main.tex
\begin{enumerate}
  \item \textbf{Independent samples.}  The $t_{m-q}$ reference distribution
  assumes independent sequenced libraries.  Repeated measures, matched samples,
  or multiple libraries from one biological unit require a subject-level
  correlation model or a blocked permutation.  FACS bins from one donor are
  a compositional partition, not four independent biological replicates; the
  GSE242880 control calibration is a diagnostic sensitivity analysis, whereas
  a purpose-built joint model such as Waterbear represents that structure
  directly.  Likewise, a 0--100\% bulk pool overlaps its sorted tails.  The
  CRISPulator analysis includes that requested design to measure its empirical
  behavior, not to redefine bulk as an independent third bin.

  \item \textbf{Enough sample-level degrees of freedom.}  Read depth cannot
  replace biological replication.  With very few samples, guide-wise
  $\widehat{\rho}$ is noisy and coefficient tests may be poorly calibrated.
  The model-based covariance treats $\widehat{\rho}$ as a fixed plug-in value,
  and the $t_{m-q}$ reference does not formally propagate its estimation
  uncertainty.  The local-equivalence result likewise requires increasing
  numbers of independent libraries; deeper sequencing at fixed replication is
  not sufficient.  The optional \texttt{bb\_calibrate\_controls()} procedure
  additionally requires all control guides to share the same residual degrees
  of freedom, as they do when every guide uses the same complete design.  Guides
  fitted with heterogeneous designs require stratified or redesigned
  calibration.  The optional \texttt{bb\_gene\_consistency()} statistic can
  recover concordant multi-guide signal when sample-level degrees of freedom
  are inadequate, but its empirical-null $p$-values measure guide consistency.
  They do not repair the missing biological replication and should not be
  described as confirmatory sample-level inference.
  Dispersion shrinkage across guides or permutation of the phenotype should be
  evaluated before treating this prototype as a production method.

  \pagebreak[2]
  \item \textbf{Mean-model form.}  A single slope assumes a linear relationship
  on the logit scale.  Residual plots should be checked.  A spline or a
  prespecified nonlinear term is preferable when the dose response is curved.
  This matters directly in viability time courses: incomplete knockout can
  produce finite depletion saturation, which Chronos models mechanistically
  but a single beta-binomial slope does not.

  \pagebreak[3]
  \item \textbf{Sparse and separated guides.}  Guides with all zero counts or
  with counts confined to one end of the phenotype range can cause boundary
  fits.  Low-count filtering should be chosen without reference to the tested
  phenotype.  A penalized or likelihood-ratio procedure is a useful future
  extension for severe separation.

  \item \textbf{Compositional counts.}  Guide counts share a fixed library
  total.  The marginal beta-binomial model handles depth and guide-specific
  overdispersion but does not remove global composition shifts.  Stable
  nontargeting controls and sensitivity analyses with alternative size factors
  are advisable.  If guides are filtered for testing or benchmarking, the
  \texttt{totals} argument must retain the unfiltered mapped-read totals.
  Recomputing totals after row subsetting changes the likelihood and can
  materially change calibration.

  \item \textbf{Guide-to-gene inference.}  The BARCS regression itself is
  guide-level and requires no Fisher or Stouffer combination.  Some replicated
  benchmark analyses use a directional Stouffer summary only as a transparent,
  explicitly labeled comparison adapter.  BARCS-GC instead estimates one
  shared guide coefficient by inverse-variance weighting and calibrates its
  Wald statistic against a robust gene-level empirical null for the
  one-replicate diagnostic.  That fixed-effect construction assumes a common
  gene effect, while its all-gene MAD assumes that fewer than half of genes are
  active.  The direction-agreement output should be inspected for guide
  heterogeneity.  We additionally evaluated BARCS-NORM, which treats the
  fitted guide coefficients within a gene as exchangeable draws
  $\widehat\beta_{gj}\sim N(\mu_g,\sigma_g^2)$ and tests the sample mean with
  a $t_{m_g-1}$ reference because $\sigma_g$ is estimated.  This literal
  normal model was calibrated but underpowered with three to six guides:
  across the 50 supporting 400-gene CRISPulator runs its average precision
  was 0.706 and its directional recall at FDR 0.10 was 0.159.  Replacing the
  $t$ reference by a
  plug-in normal tail would incorrectly treat the estimated $\sigma_g$ as
  known.  A production hierarchical gene test should therefore borrow
  variance information across genes and be benchmarked separately, ideally
  with held-out controls or phenotype permutations that preserve cross-guide
  dependence.

  \item \textbf{Processed longitudinal inputs.}  The Liang deposited values
  are normalized, batch-corrected pseudo-counts, so the fitted likelihood does
  not retain a literal raw-count interpretation.  The three-time-point
  analysis is longitudinal, but raw-read confirmation remains necessary
  before using it as a distributional validation.
  Likewise, post-fit CNV correction should be accepted only after a null-gene
  diagnostic: it succeeds in A375 and over-corrects the longitudinal HT-29
  screen.
\end{enumerate}
